% Generated draft: requires mathematical and source review.
\documentclass[11pt]{article}
\usepackage[margin=25mm]{geometry}
\usepackage{fontspec}
\setmainfont{Latin Modern Roman}
\usepackage{amsmath,amssymb,mathtools,unicode-math}
\setmathfont{Latin Modern Math}
\usepackage{enumitem,needspace,xurl,hyperref}
\setlength{\parindent}{0pt}
\setlength{\parskip}{5pt}
\setlength{\emergencystretch}{3em}
\newcommand{\N}{\mathbb{N}}
\newcommand{\Z}{\mathbb{Z}}
\newcommand{\Q}{\mathbb{Q}}
\newcommand{\R}{\mathbb{R}}
\newcommand{\C}{\mathbb{C}}
\newcommand{\F}{\mathbb{F}}
\newcommand{\A}{\mathbb{A}}
\newcommand{\PP}{\mathbb{P}}
\newcommand{\E}{\mathbb{E}}
\newcommand{\eps}{\varepsilon}
\newcommand{\dd}{\,\mathrm{d}}
\begin{document}
\section*{UnsolvedMath ID: 700977}
\textbf{Source problem code:} catalogue entry
\subsection{Definitions and mathematical statement}
A group \(G\) is \emph{residually finite} if, for every \(g\in G\setminus\{1\}\), there exist a finite group \(F\) and a homomorphism \(\varphi:G\to F\) such that \(\varphi(g)\ne 1\). For a positive integer \(d\), let
\[
\mathrm{SL}_d(\Z)=\{A\in\mathrm{Mat}_{d\times d}(\Z):\det A=1\}.
\]
An extension of shape \(\Z^d.\mathrm{SL}_d(\Z)\) is represented by a short exact sequence
\[
1\longrightarrow\Z^d\xrightarrow{\iota}G\xrightarrow{\pi}\mathrm{SL}_d(\Z)\longrightarrow1.
\]
No splitting is assumed. Conjugation on the abelian kernel induces an action \(\rho:\mathrm{SL}_d(\Z)\to\mathrm{GL}_d(\Z)\): for \(A\in\mathrm{SL}_d(\Z)\), choose any lift \(\widetilde A\in G\) and set
\[
\iota(\rho(A)v)=\widetilde A\,\iota(v)\,\widetilde A^{-1}\qquad(v\in\Z^d).
\]
This action is independent of the lift. The supplied record does not state whether arbitrary induced actions are allowed or whether the standard action \(\rho(A)v=Av\) is intended. Consequently, the group shape alone does not settle this material restriction on the question.

\textbf{Question in dimension five.} Does there exist an extension of the intended shape
\[
1\longrightarrow\Z^5\longrightarrow G\longrightarrow\mathrm{SL}_5(\Z)\longrightarrow1
\]
whose middle group \(G\) is not residually finite? Equivalently, can such an extension have an element \(g\ne1\) with \(\varphi(g)=1\) for every homomorphism \(\varphi\) from \(G\) to a finite group? The intended restriction on the induced action remains unresolved in the supplied record.

\textbf{Broader question.} The source also asks whether extensions of an arithmetic subgroup of a Chevalley group ``over a rational module'' can fail to be residually finite, apart from examples of shape \(\Z^3.\mathrm{SL}_3(\Z)\). A framework for specifying such a question is
\[
1\longrightarrow M\longrightarrow E\longrightarrow\Gamma\longrightarrow1,
\]
where \(\Gamma\) is the arithmetic subgroup and \(M\) is an additive \(\Gamma\)-module with the prescribed induced action. However, the source does not specify the Chevalley group, the arithmetic structure, the meaning of rational module, or whether the kernel is a rational vector space or an invariant integral lattice in one. It also does not specify which modules or actions are quantified over, or exactly which dimension-three examples are excluded. Thus this is a request for further counterexamples in a source-described family whose precise scope remains unresolved; the word ``usually'' supplies no mathematical quantifier.

\textbf{Supplied background qualifications.} The record asserts that every extension of the stated shape is residually finite unless possibly \(d=3\) or \(d=5\), and attributes the existence of a non-residually-finite extension of \(\Z^3\) by \(\mathrm{SL}_3(\Z)\) to the cited 1995 work of P. R. Hewitt. These are supplied background assertions, not independently verified conclusions of this draft. Their intended action restrictions require clarification. The registry labels the problem open; that label is supplied curation, not an independently verified assessment.
\subsection{Short English statement}
P. R. Hewitt asks whether an extension with kernel \(\Z^5\) and quotient \(\mathrm{SL}_5(\Z)\) can fail to be residually finite, and whether related arithmetic-group extensions provide further counterexamples beyond the cited dimension-three examples. The intended kernel actions and the scope of the broader rational-module question need clarification.
\subsection{Sources}
\begin{itemize}
\item Local UnsolvedMath registry, permanent ID \(700977\), exact source code \texttt{catalogue entry}, title \emph{Kourovka Notebook Problem 13.59}. Supplied upstream revision: \texttt{d371ee917ac9647fbb06cd6fc529deec6f3dda82}. The record was created and updated on 2026-10-16 and carries the supplied status \texttt{open}.
\item \emph{The Kourovka Notebook (Unsolved Problems in Group Theory)}, original problem 13.59, proposed by P. R. Hewitt in the 13th issue (1995). The supplied background identifies the source edition as the 21st issue edition (2026), arXiv identifier \texttt{1401.0300v48}. The supplied collection description names E. I. Khukhro and V. D. Mazurov as editors. These edition details are reproduced from the record and have not been independently checked.
\item P. R. Hewitt, contribution to \emph{Groups/St. Andrews'93 in Galway, Vol. 2}, London Mathematical Society Lecture Note Series \textbf{212}, Cambridge University Press, 1995, pp. 305--313. The source cites this contribution for the dimension-three counterexample; no contribution title or URL is supplied, and its contents have not been independently verified.
\end{itemize}
\end{document}
